\documentclass[11pt]{article}
% ============================================================================
% preamble.tex — Estilo común de los manuales de práctica SISELA (FIME / UANL)
%
% Reproduce (de forma aproximada) el estilo de los PDF originales
% (~/Descargas/manuales/Manuales/pN.pdf), de los que no existe fuente .tex.
% Compilar con:  pdflatex  (o  latexmk -pdf pN.tex)
% ============================================================================

\usepackage[utf8]{inputenc}
\usepackage[T1]{fontenc}
\usepackage[spanish,es-tabla]{babel}
\usepackage[a4paper,margin=2.5cm,headheight=15pt]{geometry}

% --- Tipografía: serif para el cuerpo, sans para los títulos ---
\usepackage{libertinus}
\usepackage[scaled=0.92]{helvet}   % sans para encabezados
\usepackage{microtype}

% --- Matemáticas ---
\usepackage{amsmath,amssymb}
\usepackage{siunitx}
\sisetup{detect-all, per-mode=symbol}

% --- Tablas y figuras ---
\usepackage{booktabs}
\usepackage{array}
\usepackage{multirow}
\usepackage{caption}
\captionsetup{font=small,labelfont=bf,labelsep=period}
\usepackage{graphicx}
\usepackage{float}

% --- Diagramas ---
\usepackage{tikz}
\usetikzlibrary{shapes.geometric,arrows.meta,positioning,calc}
\tikzset{
  block/.style   = {draw, rectangle, minimum width=2.6cm, minimum height=0.9cm,
                    align=center, font=\footnotesize},
  term/.style    = {draw, rounded corners, minimum width=2.2cm, minimum height=0.8cm,
                    align=center, font=\footnotesize, fill=black!5},
  decision/.style= {draw, diamond, aspect=2, align=center, font=\scriptsize,
                    inner sep=1pt},
  flow/.style    = {-{Latex[length=2mm]}, thick},
}

% --- Código ---
\usepackage{xcolor}
\definecolor{codebg}{gray}{0.96}
\definecolor{codekw}{rgb}{0.13,0.29,0.53}
\definecolor{codecm}{rgb}{0.40,0.45,0.40}
\usepackage{listings}
\lstdefinestyle{sisela}{
  backgroundcolor=\color{codebg}, basicstyle=\ttfamily\footnotesize,
  keywordstyle=\color{codekw}\bfseries, commentstyle=\color{codecm}\itshape,
  stringstyle=\color{codekw}, numbers=left, numberstyle=\tiny\color{gray},
  numbersep=6pt, frame=none, breaklines=true, showstringspaces=false,
  columns=fullflexible, xleftmargin=1.8em, aboveskip=1em, belowskip=1em,
  literate={á}{{\'a}}1 {é}{{\'e}}1 {í}{{\'i}}1 {ó}{{\'o}}1 {ú}{{\'u}}1
           {ñ}{{\~n}}1 {Á}{{\'A}}1 {°}{{\textdegree}}1 {µ}{{$\mu$}}1
           {✓}{{\checkmark}}1,
}
\lstset{style=sisela}
\newcommand{\code}[1]{\texttt{\small #1}}

% --- Listas ---
\usepackage{enumitem}
\setlist{nosep, leftmargin=1.6em}

% --- Encabezado / pie de página ---
\usepackage{fancyhdr}
\usepackage{lastpage}
\pagestyle{fancy}
\fancyhf{}
\fancyhead[L]{\footnotesize\sffamily Sistemas Eléctricos de Aeronaves — Práctica \practnum}
\fancyhead[R]{\footnotesize\sffamily FIME / UANL}
\fancyfoot[C]{\footnotesize P\'agina \thepage\ de \pageref{LastPage}}
\renewcommand{\headrulewidth}{0.4pt}
\renewcommand{\footrulewidth}{0.4pt}

% --- Títulos de sección en sans-serif ---
\usepackage{titlesec}
\titleformat{\section}{\normalfont\Large\sffamily\bfseries}{\thesection}{0.7em}{}
\titleformat{\subsection}{\normalfont\large\sffamily\bfseries}{\thesubsection}{0.6em}{}
\titleformat{\subsubsection}{\normalfont\normalsize\sffamily\bfseries}{\thesubsubsection}{0.5em}{}

% --- hyperref al final ---
\usepackage[hidelinks,pdfusetitle]{hyperref}

% ============================================================================
% Comandos de portada y cajas
% ============================================================================
\newcommand{\practnum}{N}          % lo redefine cada pN.tex
\newcommand{\practtitle}{}         % idem

\newcommand{\portada}{%
  \thispagestyle{empty}
  \begin{center}
    {\large Universidad Autónoma de Nuevo León}\\[2pt]
    {\normalsize Facultad de Ingeniería Mecánica y Eléctrica}\\[2.2cm]
    \rule{\linewidth}{0.4pt}\\[0.4cm]
    {\Huge\sffamily\bfseries Práctica \practnum}\\[0.5cm]
    {\large \practtitle}\\[0.3cm]
    \rule{\linewidth}{0.4pt}\\[2.5cm]
    {\normalsize Laboratorio de Sistemas Electrónicos de Aeronaves\\(SISELA)}\\[1.5cm]
    {\normalsize Ciclo Escolar 2026}\\[3cm]
  \end{center}
}

\newenvironment{resumen}{\par\noindent\rule{\linewidth}{0.4pt}\par\smallskip
  \small\itshape\noindent\textbf{Resumen Ejecutivo} ---\ }{\par\smallskip
  \noindent\rule{\linewidth}{0.4pt}\par}

\newenvironment{nota}[1][Nota]{\par\smallskip\noindent\rule{\linewidth}{0.4pt}\par\smallskip
  \small\itshape\noindent\textbf{#1} ---\ }{\par\smallskip
  \noindent\rule{\linewidth}{0.4pt}\par\smallskip}

\renewcommand{\practnum}{6}
\renewcommand{\practtitle}{Control de Actuadores II: Conmutación de Potencia con BJT\\
Interfaces de Alta Corriente — Énfasis Aeronáutico}

\begin{document}
\portada

\begin{resumen}
Esta sesión se enfoca en el control de dispositivos que demandan mucha más potencia de
la que un microcontrolador puede entregar directamente: motores, solenoides, luces de
alta intensidad. El transistor de unión bipolar (BJT), actuando como interruptor
electrónico, permite que una señal de baja potencia del GPIO controle el flujo de una
corriente mucho mayor hacia la carga. La práctica cubre tres casos de estudio
progresivos: conmutación de un \emph{buzzer} con el 2N2222A (señal), control de un motor
DC de 12\,V con el TIP31C mediante un driver NPN (\emph{low-side}), y conmutación
\emph{high-side} con el TIP32C. Para cada caso se presentan los cálculos de diseño
detallados, se verifica mediante simulación en Multisim y se implementa
experimentalmente con firmware MicroPython y C++.
\end{resumen}

\tableofcontents
\newpage

% ============================================================================
\section{Introducción}

Los microcontroladores (MCUs) como el ESP32 y el RP2040 son el cerebro de los sistemas
embebidos aeronáuticos, operando con voltajes y corrientes bajos (típicamente 3.3\,V y
algunos miliamperios por pin GPIO). Sin embargo, el mundo real — especialmente en
aplicaciones aeronáuticas — requiere controlar dispositivos que demandan mucha más
potencia: motores para bombas de combustible o ventilación, solenoides para válvulas
hidráulicas o neumáticas, luces de indicación o iluminación de alta intensidad, y
mecanismos de actuación. Aquí es donde los \textbf{transistores}, actuando como
\textbf{interruptores electrónicos}, se vuelven indispensables. Permiten que una señal
de baja potencia del MCU controle el flujo de una corriente mucho mayor hacia la carga.

En la Práctica 5 se introdujo el control de posición con servomotores, donde el PWM
codificaba un ángulo. En esta práctica se da el siguiente paso: utilizar PWM para regular
la \textbf{potencia entregada} a cargas de mayor demanda energética. El transistor de
unión bipolar (BJT) opera como un interruptor controlado por corriente: una pequeña
corriente de base (proveniente del GPIO a través de una resistencia calculada) permite o
impide el paso de una corriente de colector mucho mayor hacia la carga.

El diseño riguroso de estas interfaces de potencia es un aspecto crucial en los sistemas
electrónicos aeronáuticos. Un transistor que no alcanza la saturación disipará calor
excesivo, reduciendo la fiabilidad; la ausencia de protección contra transitorios
inductivos puede destruir el silicio; y una selección inadecuada del esquema de
conmutación (\emph{low-side} vs \emph{high-side}) puede comprometer la seguridad ante
fallas de aislamiento.

Esta práctica cubre tres casos de estudio progresivos: conmutación de un \emph{buzzer}
con el 2N2222A (señal), control de un motor DC de 12\,V con el TIP31C mediante un driver
NPN, y conmutación \emph{high-side} con el TIP32C. Para cada caso se presentan los
cálculos de diseño detallados, se verificará mediante simulación en Multisim y se
implementará experimentalmente con firmware MicroPython y C++.

% ============================================================================
\section{Objetivos de Aprendizaje}
\begin{enumerate}
  \item Explicar el funcionamiento del BJT como interruptor: regiones de corte y
        saturación, y los parámetros clave del \emph{datasheet}
        ($h_{FE}$, $V_{BE(\text{sat})}$, $P_{D,\max}$, $R_{\theta JA}$).
  \item Diseñar circuitos de conmutación BJT seleccionando el transistor adecuado
        (NPN/PNP, señal/potencia) y la configuración (\emph{low-side}/\emph{high-side})
        según la carga y la tensión de alimentación.
  \item Calcular la corriente de base de diseño aplicando un factor de sobre-excitación
        (\emph{overdrive factor}, ODF) para asegurar la saturación:
        $I_{B,\text{design}} = (I_C/h_{FE,\min})\times\text{ODF}$.
  \item Determinar el valor y la potencia de las resistencias de base ($R_B$) para el
        transistor de potencia y su driver, verificando que no se excedan las
        especificaciones nominales.
  \item Analizar la disipación de potencia en el transistor
        ($P_D = V_{CE(\text{sat})}\times I_C$) y evaluar la necesidad de disipador
        consultando $R_{\theta JA}$ y $T_{J,\max}$ del \emph{datasheet}.
  \item Seleccionar e implementar el diodo de protección (\emph{flyback}) para cargas
        inductivas, considerando $I_{F(\text{AV})}\ge I_\text{carga}$ y
        $V_{RRM} > V_{CC}$.
  \item Diseñar la etapa de \emph{driver} NPN para controlar un transistor PNP de
        potencia desde un MCU de 3.3\,V, incluyendo los cálculos de la resistencia de
        base del PNP.
  \item Simular los circuitos en Multisim y verificar corrientes ($I_C$, $I_B$),
        voltajes ($V_{CE}$, $V_{BE}$) y la operación del diodo \emph{flyback} mediante
        análisis transitorio.
  \item Implementar y verificar experimentalmente montando los circuitos, midiendo
        $V_{CE(\text{sat})}$ y controlando la carga con firmware de 4 modos
        (on/off, PWM manual, barrido, potenciómetro).
  \item Contextualizar las interfaces de potencia en aplicaciones aeronáuticas: control
        de luces, bombas, solenoides, comandos discretos ARINC 429 y prácticas de
        conexión a tierra (\emph{grounding}).
\end{enumerate}

% ============================================================================
\section{Fundamentos Teóricos}

\subsection{El transistor BJT como interruptor}
En sistemas de aeronaves, el microcontrolador no puede energizar directamente una bomba
de combustible o un solenoide de tren de aterrizaje, ya que estas cargas requieren
corrientes de cientos de miliamperios o amperios. El transistor BJT se utiliza como
interruptor electrónico operando en dos regiones:
\begin{itemize}
  \item \textbf{Corte}: $I_B=0 \Rightarrow$ el transistor es un circuito abierto
        ($I_C\approx 0$, $V_{CE}\approx V_{CC}$).
  \item \textbf{Saturación}: $I_B > I_C/h_{FE,\min} \Rightarrow$ el transistor es un
        interruptor cerrado ($V_{CE(\text{sat})}$ mínima, típicamente
        \SIrange{0.2}{1.4}{\volt} según el dispositivo).
\end{itemize}
La región activa (amplificación lineal) \textbf{debe evitarse} en conmutación, ya que el
transistor disipa potencia excesiva ($P=V_{CE}\times I_C$, con $V_{CE}$ intermedio).

\begin{nota}[Región activa vs saturación]
En amplificación, se busca la región activa. En conmutación de potencia, se busca pasar
rápidamente de corte a saturación y viceversa, minimizando el tiempo en la región lineal
para reducir pérdidas de conmutación.
\end{nota}

\subsubsection{Factor de sobre-excitación (\emph{Overdrive Factor}, ODF)}
Para garantizar que el transistor permanezca en saturación bajo todas las condiciones
(temperatura, tolerancias, degradación), se diseña la corriente de base con un margen de
seguridad:
\begin{align}
  I_{B,\min} &= \frac{I_C}{h_{FE,\min}}, \\
  I_{B,\text{design}} &= I_{B,\min}\times\text{ODF}.
\end{align}
Valores típicos de ODF: 2--5 para transistores de señal, 3--10 para transistores de
potencia (donde $h_{FE}$ es bajo y variable con la temperatura).

\begin{table}[H]\centering
\caption{Parámetros críticos del \emph{datasheet} para diseño de conmutación.}
\begin{tabular}{ll}
\toprule
Parámetro & Descripción \\
\midrule
$h_{FE,\min}$ & Ganancia de corriente DC en saturación (varía con $I_C$ y $T$) \\
$V_{BE(\text{sat})}$ & Caída base--emisor en saturación (0.7--1.8\,V según $I_C$) \\
$V_{CE(\text{sat})}$ & Caída colector--emisor en saturación (0.2--1.4\,V) \\
$I_{C,\max}$ & Corriente de colector continua máxima \\
$P_{D,\max}$ & Disipación máxima sin disipador (a 25\si{\celsius}) \\
$V_{CEO}$ & Máximo voltaje C--E con base abierta \\
$R_{\theta JA}$ & Resistencia térmica juntura--ambiente (\si{\celsius\per\watt}) \\
\bottomrule
\end{tabular}
\end{table}

\subsection{Cálculo de la resistencia de base ($R_B$)}
La resistencia de base limita la corriente de base al valor de diseño:
\begin{equation}
  R_B = \frac{V_{OH} - V_{BE(\text{sat})}}{I_{B,\text{design}}},
\end{equation}
donde $V_{OH}$ es el voltaje de salida alto del GPIO (3.3\,V para ESP32/RP2040). La
potencia disipada en $R_B$ es
\begin{equation}
  P_{R_B} = I_{B,\text{design}}^2\times R_B = (V_{OH}-V_{BE(\text{sat})})\times I_{B,\text{design}}.
\end{equation}

\subsection{Configuración \emph{low-side} (NPN)}
En la configuración \emph{low-side}, el transistor NPN se sitúa entre la carga y tierra.
La carga se conecta desde la fuente de alimentación ($V_{CC,\text{Load}}$) al colector
del transistor.

\begin{figure}[H]\centering
\begin{tikzpicture}[node distance=1.1cm]
  \node[term] (vcc) {$+V_{CC}$};
  \node[block, below=1.4cm of vcc, xshift=-1.6cm] (load) {Carga};
  \node[block, below=1.4cm of vcc, xshift=1.6cm] (fb) {\emph{Flyback}};
  \node[block, below=2.6cm of load] (npn) {NPN};
  \node[term, left=1.8cm of npn] (gpio) {GPIO};
  \node[term, below=1.2cm of npn] (gnd) {GND};
  \draw[flow] (vcc) -| (load); \draw[flow] (vcc) -| (fb);
  \draw[flow] (load) -- (npn); \draw[flow] (fb) -- (npn);
  \draw[flow] (gpio) -- node[above,font=\scriptsize]{$R_B$} (npn);
  \draw[flow] (npn) -- (gnd);
\end{tikzpicture}
\caption{Configuración \emph{low-side} con transistor NPN y diodo \emph{flyback}.}
\end{figure}

\textbf{Operación}: GPIO HIGH $\Rightarrow I_B$ fluye $\Rightarrow$ NPN satura
$\Rightarrow$ corriente fluye por la carga. GPIO LOW $\Rightarrow I_B=0
\Rightarrow$ NPN en corte $\Rightarrow$ carga apagada.

\subsection{Configuración \emph{high-side} (PNP con driver NPN)}
En aplicaciones aeronáuticas se prefiere la conmutación \emph{high-side} (cortar la línea
positiva) porque si un cable de señal toca accidentalmente la estructura metálica del
avión (tierra), la carga \textbf{no} se energiza de forma involuntaria. Para controlar un
PNP de potencia desde un MCU de 3.3\,V que conmuta una carga de 12\,V, se requiere un
transistor NPN como \emph{driver} de nivel.

\textbf{Operación}: GPIO HIGH $\Rightarrow$ Q1 (driver NPN) satura $\Rightarrow$ la base
de Q2 (PNP) se lleva a $\sim V_{CE(\text{sat}),Q1} \Rightarrow$ Q2 satura $\Rightarrow$
carga ON. GPIO LOW $\Rightarrow$ Q1 en corte $\Rightarrow$ base de Q2 sin camino a
tierra $\Rightarrow$ Q2 en corte $\Rightarrow$ carga OFF.

\subsection{Diodo de protección (\emph{Flyback})}
Las cargas inductivas (motores, solenoides, relés, bobinas) almacenan energía en su campo
magnético. Al apagar el transistor abruptamente, la corriente que circulaba por la
inductancia intenta mantenerse según la ley de Lenz, generando un pico de voltaje inverso
que puede alcanzar cientos de volts y destruir el transistor. El diodo \emph{flyback} se
coloca en paralelo con la carga (cátodo al positivo de la alimentación, ánodo al colector
del NPN) para proporcionar un camino de recirculación para la corriente inductiva.

\begin{table}[H]\centering
\caption{Criterios de selección del diodo \emph{flyback}.}
\begin{tabular}{lll}
\toprule
Parámetro & Criterio & Observación \\
\midrule
$I_{F(\text{AV})}$ & $\ge I_\text{carga}$ & Corriente directa promedio \\
$V_{RRM}$ & $> V_{CC,\text{Load}}$ & Voltaje inverso repetitivo máximo \\
$t_{rr}$ & Lo más bajo posible & Schottky preferidos ($t_{rr}\approx 0$) \\
Tipo & 1N5817--1N5819 (Schottky) & $V_F$ bajo (0.3\,V vs 0.7\,V), recuperación rápida \\
     & 1N4001--1N4007 (rectificador) & Económico, $V_{RRM}$ alto, $t_{rr}$ más lento \\
\bottomrule
\end{tabular}
\end{table}

\subsection{Disipación de potencia y gestión térmica}
Un transistor en saturación no es un conductor perfecto: la caída $V_{CE(\text{sat})}$
genera pérdidas:
\begin{equation}
  P_D = V_{CE(\text{sat})}\times I_C.
\end{equation}
La temperatura de la juntura se estima como
\begin{equation}
  T_J = T_A + P_D\times R_{\theta JA}.
\end{equation}
Si $T_J > T_{J,\max}$ (típicamente \SI{150}{\celsius}), se requiere un disipador de calor
(\emph{heatsink}) que reduce la resistencia térmica efectiva.

\begin{nota}[Entorno aeronáutico]
En aviónica, la temperatura ambiente dentro de un \emph{bay} de equipos puede alcanzar
\SIrange{55}{70}{\celsius}. El diseño de gestión térmica debe considerar este margen
reducido, además de restricciones de espacio, peso y vibración para la sujeción del
disipador.
\end{nota}

\subsection{Control discreto vía ARINC 429}
La activación de bombas, luces o solenoides en aeronaves modernas se comanda mediante
palabras ARINC 429 de tipo \textbf{discreto}. A diferencia del formato BNR (datos
continuos como la altitud o el ángulo de alerón), las palabras discretas asignan bits
específicos del campo de datos (bits 11--29) a estados ON/OFF de sistemas individuales.

\begin{table}[H]\centering
\caption{Estructura de una palabra ARINC 429 discreta (ejemplo Label 270).}
\begin{tabular}{llll}
\toprule
Bits & Campo & Contenido & Ejemplo \\
\midrule
1--8 & Label & 270 (octal, invertido) & 10111000 \\
9--10 & SDI & Source ID & 00 \\
11 & Bit 11 & Bomba combustible aux.\ (1 = ON) & \\
12 & Bit 12 & \emph{Master Caution} (1 = Activar) & \\
13--29 & Bits 13--29 & Otros discretos & --- \\
29--30 & SSM & Status (00 = Normal) & 00 \\
31 & P & Paridad impar & Calculada \\
\bottomrule
\end{tabular}
\end{table}
El sistema de potencia debe validar la paridad y el SSM antes de actuar sobre los
transistores de salida, evitando activaciones espurias por ruido en el bus.

% ============================================================================
\section{Consideraciones de Seguridad y Buenas Prácticas}
\begin{enumerate}
  \item \textbf{Transitorios inductivos}: las bobinas de motores y solenoides almacenan
        energía magnética. Al apagar el transistor, esta energía genera un arco de
        voltaje inverso que puede destruir el silicio. El uso de diodos \emph{flyback}
        es \textbf{obligatorio} para cargas inductivas.
  \item \textbf{Disipación de calor}: verificar siempre que $T_J < T_{J,\max}$
        (ecuación 6). Si se excede, instalar disipador. Monitorizar la temperatura del
        encapsulado durante las pruebas.
  \item \textbf{Tierra común (GND)}: la tierra del MCU y la tierra de la fuente de
        potencia \textbf{deben estar conectadas}. Sin GND común, la referencia de
        voltaje del GPIO es indefinida respecto a la base del transistor.
  \item \textbf{Aislamiento de buses}: mantener el bus de 12\,V de potencia físicamente
        separado de las líneas de 3.3\,V de datos para evitar ruidos de conmutación en
        el ADC.
  \item \textbf{No alimentar cargas desde 3V3}: el regulador 3.3\,V del MCU tiene límite
        de corriente ($\sim$\SI{500}{\milli\ampere}). Usar siempre fuente externa para
        la carga.
  \item \textbf{Polaridad de los transistores}: verificar configuración de pines (B, C,
        E) según el encapsulado (TO-92 para 2N2222A, TO-220 para TIP31C/TIP32C). Una
        conexión invertida puede destruir el dispositivo.
\end{enumerate}

% ============================================================================
\section{Equipo y Materiales}
\begin{table}[H]\centering
\caption{Lista de componentes para la práctica.}
\begin{tabular}{cll l}
\toprule
Cant.\ & Componente & Especificación & Uso \\
\midrule
1 & Microcontrolador & ESP32 DevKit / RP2040 & Control GPIO/PWM \\
1 & Cable USB & Micro-B / USB-C & Programación y alimentación \\
1 & Protoboard & 830 puntos & Prototipado \\
--- & Alambre AWG 22 & Varios colores & Conexiones \\
1 & 2N2222A & NPN, TO-92 & Señal (buzzer) / \emph{Driver} \\
1 & TIP31C & NPN potencia, TO-220 & Motor DC (\emph{low-side}) \\
1 & TIP32C & PNP potencia, TO-220 & Motor DC (\emph{high-side}) \\
1 & Buzzer & 5\,V, 20\,mA & Carga Caso 1 \\
1 & Motor DC & 12\,V, 0.5\,A nom. / 1.5\,A \emph{stall} & Cargas Casos 2--3 \\
--- & Resistencias & Varios valores, 1/4\,W y 5\,W & $R_B$ y $R_\text{pullup}$ \\
1 & 1N4001 (o 1N4007) & Rectificador & \emph{Flyback} (buzzer) \\
1 & 1N5817 (o 1N5819) & Schottky & \emph{Flyback} (motor) \\
1 & Condensador & \SI{100}{\micro\farad} / \SI{16}{\volt} & Desacoplo fuente \\
1 & Disipador TO-220 & Opcional & TIP31C / TIP32C \\
1 & Potenciómetro & \SI{10}{\kilo\ohm} & Modo 4 (ADC) \\
\bottomrule
\end{tabular}
\end{table}

\paragraph{Equipo de laboratorio.} Laptop con IDE (Thonny/VS Code + Pymakr) y monitor
serie; fuente de alimentación DC variable/dual (3.3\,V, 5\,V, 12\,V con $\ge$\,2\,A);
multímetro digital (DMM) para medir $V_{CE(\text{sat})}$, $I_B$, $I_C$; osciloscopio
(recomendado) para verificar PWM y transitorios; software de simulación (NI Multisim o
equivalente).

% ============================================================================
\section{Etapa de Diseño}

\subsection{Caso de estudio 1: buzzer con 2N2222A (\emph{low-side})}
Conmutar un buzzer de 3.3\,V / 20\,mA usando un transistor NPN de señal (2N2222A) en
configuración \emph{low-side}. El circuito funciona como \textbf{anunciador acústico de
\emph{warning}} (alerta de sobretemperatura o tren de aterrizaje).

Datos: $I_C=\SI{20}{\milli\ampere}$; \emph{datasheet} 2N2222A: $h_{FE,\min}=50$,
$V_{BE(\text{sat}),\max}=\SI{1.0}{\volt}$, $V_{CE(\text{sat})}\approx\SI{0.2}{\volt}$.
\begin{align}
  I_{B,\min} &= \frac{I_C}{h_{FE,\min}} = \frac{0.020}{50} = \SI{0.4}{\milli\ampere}, \\
  I_{B,\text{design}} &= I_{B,\min}\times\text{ODF} = 0.4\times 5 = \SI{2.0}{\milli\ampere}
        \quad (\text{ODF}=5), \\
  R_B &= \frac{V_{OH}-V_{BE(\text{sat})}}{I_{B,\text{design}}} = \frac{3.3-1.0}{0.002}
        = \SI{1150}{\ohm}.
\end{align}
Valor estándar: $R_B=\SI{1.2}{\kilo\ohm}$. Potencia disipada en el transistor:
$P_D = V_{CE(\text{sat})}\times I_C = 0.2\times 0.020 = \SI{4}{\milli\watt}$.
\textbf{No requiere disipador} ($P_{D,\max,\text{2N2222A}}=\SI{625}{\milli\watt}$).

\subsection{Caso de estudio 2: motor DC 12\,V con TIP31C (\emph{low-side})}
Conmutar un motor DC de 12\,V (1.5\,A en \emph{stall}) usando el TIP31C como transistor
de potencia, con el 2N2222A como \emph{driver} inversor. Simula el control de una
\textbf{bomba auxiliar de combustible}. El TIP31C tiene $h_{FE,\min}=10$ a
$I_C=\SI{1.5}{\ampere}$, lo que requiere $I_B\approx$ \SIrange{150}{450}{\milli\ampere}.
Un GPIO de 3.3\,V no puede suministrar esta corriente directamente, por lo que se utiliza
una etapa de \emph{driver} con \emph{pull-up} a 12\,V.

\begin{nota}[Lógica invertida]
En esta configuración la lógica está invertida: GPIO HIGH $\Rightarrow$ Q1 ON
$\Rightarrow$ base de Q2 baja $\Rightarrow$ Q2 OFF $\Rightarrow$ motor apagado. GPIO LOW
$\Rightarrow$ Q1 OFF $\Rightarrow$ $R_\text{PU}$ suministra $I_B$ a Q2 $\Rightarrow$ Q2 ON
$\Rightarrow$ motor encendido. El firmware debe invertir la lógica en software.
\end{nota}

\textbf{Cálculos para TIP31C (Q2)}: $I_{C,Q2}=I_\text{load,max}=\SI{1.5}{\ampere}$
(\emph{stall}); $h_{FE,\min,Q2}=10$, $V_{BE(\text{sat}),Q2,\max}=\SI{1.8}{\volt}$,
$V_{CE(\text{sat}),Q2,\max}=\SI{1.4}{\volt}$.
\begin{align}
  I_{B,Q2,\min} &= \frac{1.5}{10} = \SI{150}{\milli\ampere}, &
  I_{B,Q2,\text{design}} &= 150\times 3 = \SI{450}{\milli\ampere}\ (\text{ODF}=3).
\end{align}

\textbf{Resistencia de \emph{pull-up} $R_\text{PU}$} (suministra $I_{B,Q2}$ cuando Q1 está
OFF):
\begin{equation}
  R_\text{PU} = \frac{V_{CC}-V_{BE(\text{sat}),Q2}}{I_{B,Q2,\text{design}}}
              = \frac{12-1.8}{0.45} = \SI{22.67}{\ohm}\approx\SI{22}{\ohm}.
\end{equation}
Potencia: $P_{R_\text{PU}} = I_{B,Q2}^2\times R_\text{PU} = 0.45^2\times 22 = \SI{4.46}{\watt}$
— se requiere resistencia de potencia $\ge\SI{5}{\watt}$.

\textbf{Cálculos para 2N2222A (Q1, \emph{driver})}: cuando Q1 está ON, desvía la
corriente de $R_\text{PU}$: $I_{C,Q1} = (V_{CC}-V_{CE(\text{sat}),Q1})/R_\text{PU} =
(12-0.3)/22 = \SI{532}{\milli\ampere}$ (aceptable, $I_{C,\max}\approx$
\SIrange{600}{800}{\milli\ampere}). Con $h_{FE,\min,Q1}\approx 35$ a
$I_C=\SI{532}{\milli\ampere}$: $I_{B,Q1,\min}=532/35=\SI{15.2}{\milli\ampere}$,
$I_{B,Q1,\text{design}}=15.2\times 1.5=\SI{22.8}{\milli\ampere}$ (ODF=1.5).
\begin{equation}
  R_{B1} = \frac{3.3-1.2}{0.0228} = \SI{92.1}{\ohm}\approx\SI{100}{\ohm}.
\end{equation}

\textbf{Análisis térmico de TIP31C (Q2)}:
$P_{D,Q2}=V_{CE(\text{sat}),Q2}\times I_{C,Q2}=1.4\times 1.5=\SI{2.1}{\watt}$;
$T_J=T_A+P_D\times R_{\theta JA}=25+2.1\times 62.5=\SI{156}{\celsius} >
T_{J,\max}=\SI{150}{\celsius} \Rightarrow$ \textbf{se requiere disipador}. Con
$R_{\theta HA}=\SI{15}{\celsius\per\watt}$:
$T_J=25+2.1\times(1+0.5+15)=\SI{59.7}{\celsius}$ \checkmark.

\subsection{Caso de estudio 3: motor DC 12\,V con TIP32C (\emph{high-side})}
Conmutar el mismo motor DC de 12\,V con conmutación \emph{high-side} usando el TIP32C
(PNP potencia) controlado por el 2N2222A (\emph{driver} NPN). Esta configuración es la
preferida en aeronáutica por las razones de seguridad descritas en la \S3.4.

\textbf{Cálculos para TIP32C (Q2)}: $I_{C,Q2}=\SI{1.5}{\ampere}$, $h_{FE,\min,Q2}=10$,
$|V_{EB(\text{sat}),Q2,\max}|=\SI{1.8}{\volt}$, $|V_{EC(\text{sat}),Q2,\max}|=\SI{1.4}{\volt}$.
\begin{align}
  I_{B,Q2,\min} &= \frac{1.5}{10} = \SI{150}{\milli\ampere}, &
  I_{B,Q2,\text{design}} &= 150\times 2 = \SI{300}{\milli\ampere}\ (\text{ODF}=2).
\end{align}
\textbf{Cálculos para 2N2222A (Q1, \emph{driver})}: cuando Q1 está ON, debe hundir
(\emph{sink}) $I_{C,Q1}=I_{B,Q2,\text{design}}=\SI{300}{\milli\ampere}$; con
$h_{FE,\min,Q1}\approx 35$: $I_{B,Q1,\min}=300/35=\SI{8.6}{\milli\ampere}$,
$I_{B,Q1,\text{design}}=8.6\times 1.5=\SI{12.9}{\milli\ampere}$ (ODF=1.5).
\begin{equation}
  R_{B1} = \frac{3.3-1.0}{0.0129} = \SI{178}{\ohm}\approx\SI{180}{\ohm}.
\end{equation}
$R_{B2}$ (base de Q2, conecta el colector de Q1 a la base de Q2): cuando Q1 ON,
$V_{C,Q1}=V_{CE(\text{sat}),Q1}\approx\SI{0.3}{\volt}$; la base de Q2 está a
$V_{B,Q2}=V_{E,Q2}-|V_{EB(\text{sat})}|=12-1.8=\SI{10.2}{\volt}$:
$V_{R_{B2}}=10.2-0.3=\SI{9.9}{\volt}$, $R_{B2}=9.9/0.3=\SI{33}{\ohm}$.
Potencia en $R_{B2}$: $(0.3)^2\times 33=\SI{2.97}{\watt}\Rightarrow$ resistencia
$\ge\SI{5}{\watt}$.

\textbf{Análisis térmico}: TIP32C (Q2) — $P_D=|V_{EC(\text{sat})}|\times I_C=1.4\times
1.5=\SI{2.1}{\watt}$, \textbf{requiere disipador} (mismo análisis que el Caso 2). 2N2222A
(Q1) — $P_D=0.3\times 0.3=\SI{0.09}{\watt}$, OK sin disipador.

\begin{table}[H]\centering
\caption{Comparación de los tres casos de estudio.}
\begin{tabular}{lccc}
\toprule
Parámetro & Caso 1 (Buzzer) & Caso 2 (Motor, \emph{low-side}) & Caso 3 (Motor, \emph{high-side}) \\
\midrule
Transistor de potencia & 2N2222A & TIP31C & TIP32C (PNP) \\
\emph{Driver} & Directo (GPIO) & 2N2222A + $R_\text{PU}$ & 2N2222A + $R_{B2}$ \\
$I_C$ (máx) & 20\,mA & 1.5\,A & 1.5\,A \\
$R_B$ principal & 1.2\,k$\Omega$ (1/4\,W) & 22\,$\Omega$ (5\,W) & 33\,$\Omega$ (5\,W) \\
$P_D$ transistor & 4\,mW & 2.1\,W & 2.1\,W \\
Disipador & No & Sí & Sí \\
Diodo \emph{flyback} & 1N4001 & 1N5819 & 1N5819 \\
Lógica & Directa & Invertida & Directa \\
Ventaja aeronáutica & Simplicidad & Alta corriente & Seguridad ante fallas \\
\bottomrule
\end{tabular}
\end{table}

\subsection{Asignación de pines por plataforma}
\begin{table}[H]\centering
\caption{Asignación de pines según plataforma.}
\begin{tabular}{llll}
\toprule
Señal & ESP32 & RP2040 & Descripción \\
\midrule
PWM / GPIO & GPIO18 & GP18 (PWM1A) & Señal a $R_B$ del transistor \\
GND & GND & GND & Tierra común \\
ADC (modo 4) & GPIO34 & GP26 (ADC0) & Potenciómetro \\
VCC carga & Fuente ext.\ 12\,V & Fuente ext.\ 12\,V & Alimentación del motor \\
\bottomrule
\end{tabular}
\end{table}

\begin{nota}[Frecuencia PWM]
A diferencia de la P5 (servo a 50\,Hz), en esta práctica la frecuencia PWM es 1\,kHz
(configurable). Valores típicos según la carga: LED/lámpara 100--1000\,Hz, motor DC
500--2000\,Hz, electroválvula 50--200\,Hz.
\end{nota}

% ============================================================================
\section{Procedimiento Experimental}

\subsection{Paso 1: preparación del entorno}
\begin{enumerate}
  \item Instalar MicroPython en la placa (o configurar PlatformIO con
        \code{PRACTICE=6}).
  \item Verificar que todos los componentes estén disponibles.
  \item Con la fuente apagada, montar el circuito del Caso 1 (buzzer + 2N2222A).
  \item Verificar con multímetro: continuidad de GND común, polaridad correcta.
  \item Sincronizar archivos del repositorio (\code{main.py}, \code{boot.py},
        \code{lib/actuator\_pwm.py}).
  \item Al reiniciar, debe aparecer: \code{=== Práctica 6 --- Conmutación PWM ===}.
\end{enumerate}

\subsection{Modos de operación del firmware}
\begin{table}[H]\centering
\caption{Modos de operación del firmware P6.}
\begin{tabular}{cll}
\toprule
Modo & Nombre & Descripción \\
\midrule
1 & Encendido/Apagado & Alterna \emph{duty} 0\,\%/100\,\% cada segundo \\
2 & PWM manual & Ingresa \emph{duty} (0--100) por terminal \\
3 & Barrido automático & Rampa 0\,\%$\to$100\,\%$\to$0\,\% continua \\
4 & Potenciómetro (ADC) & Lectura ADC mapea a \emph{duty} \\
\bottomrule
\end{tabular}
\end{table}
Tecla \code{m} $\to$ regresar al menú en cualquier modo.

\subsection{Caso 1: anunciador acústico (\emph{Warning})}
Ejecutar Modo 1. \textbf{Verificar}: el buzzer suena con \emph{duty} 100\,\% y se
silencia con 0\,\%. Medir $V_{CE(\text{sat})}$ del 2N2222A (esperado $<\SI{0.3}{\volt}$).
Si se dispone de osciloscopio, observar la señal PWM en la base del transistor.

\subsection{Caso 2: bomba auxiliar (motor DC 12\,V, \emph{low-side})}
Con la fuente apagada, montar el circuito del Caso 2. Usar resistencia de potencia
($\ge\SI{5}{\watt}$) para $R_\text{PU}$. Instalar disipador en el TIP31C. Conectar GND
del MCU al GND de la fuente de 12\,V. Ejecutar Modo 1 (recordar la lógica invertida).
Verificar giro del motor; medir $V_{CE(\text{sat}),Q2}$ (esperado $<\SI{1.4}{\volt}$;
ideal $<\SI{0.5}{\volt}$). Alternar con Modo 2: variar el \emph{duty} de 0 a 100\,\% y
observar la velocidad del motor.

\begin{nota}[Transitorio del diodo]
Si se dispone de osciloscopio, capturar la forma de onda en el colector de Q2 al momento
de apagar el motor (\emph{duty} pasa de 100\,\% a 0\,\%). Sin diodo \emph{flyback}, se
observará un pico de voltaje que puede superar $V_{CEO}$ del TIP31C (100\,V). Con diodo,
el pico se limita a $\sim V_{CC}+V_F\approx\SI{12.3}{\volt}$.
\end{nota}

\subsection{Caso 3: conmutación \emph{high-side} (TIP32C)}
Reconfigurar al circuito del Caso 3. Usar resistencia de potencia ($\ge\SI{5}{\watt}$)
para $R_{B2}$. Instalar disipador en el TIP32C. Verificar conexiones: emisor de Q2 a
+12\,V, colector de Q2 al motor, motor a GND. Ejecutar Modo 1: la lógica es
\textbf{directa} en esta configuración (GPIO HIGH $\Rightarrow$ motor ON). Medir
$|V_{EC(\text{sat}),Q2}|$ (esperado $<\SI{1.4}{\volt}$). Comparar cualitativamente la
operación con el Caso 2.

\subsection{Caso 4: control por palabra discreta (ARINC 429 simulado)}
Este caso integra el control de potencia con un protocolo de aviónica. El firmware recibe
una trama de 32 bits por serie (en hexadecimal) y actúa sobre los transistores según los
bits discretos: valida la paridad impar, verifica el \emph{Label} (270 octal), extrae el
bit 11 (bomba combustible) y el bit 12 (\emph{Master Caution}), y comanda los pines
correspondientes.

\begin{figure}[H]\centering
\begin{tikzpicture}[node distance=1.0cm]
  \node[term] (a) {Trama UART (hex)};
  \node[block, below=of a] (b) {Validar paridad\\e integridad};
  \node[decision, below=1.4cm of b] (c) {¿Bit 11\\== 1?};
  \node[block, right=2.2cm of c] (d) {Desactivar\\transistor};
  \node[block, below=1.4cm of c] (e) {Saturar transistor\\(bomba ON)};
  \draw[flow] (a)--(b); \draw[flow] (b)--(c);
  \draw[flow] (c) -- node[above,font=\scriptsize]{No} (d);
  \draw[flow] (c) -- node[right,font=\scriptsize]{Sí} (e);
\end{tikzpicture}
\caption{Lógica de control de potencia mediante comandos discretos ARINC 429.}
\end{figure}

\subsection{Implementación C++ (PlatformIO)}
El archivo \code{p6.cpp} implementa tres modos usando la clase \code{PropulsionSystem}
(\code{common/propulsion.h}, API en porcentaje 0--100). Compilar con
\code{pio run -e esp32dev} o \code{pio run -e pico} (con \code{-DPRACTICE=6}).

% ============================================================================
\section{Alternativa de Trabajo en Casa (Multímetro + Simulación)}
\label{sec:casa}

\paragraph{Con multímetro (en casa).} En modo DC, medir $V_{CE(\text{sat})}$ directamente
entre colector y emisor mientras el transistor está saturado (Modo 1 sostenido en ON) —
comparar contra los valores esperados de la Tabla de comparación. Insertar el multímetro
en modo amperímetro (en serie, con precaución de rango) para medir $I_B$ e $I_C$ y
verificar el cálculo de $R_B$. Usar el modo diodo del multímetro para comprobar la
polaridad y el estado del diodo \emph{flyback} antes de energizar. Si el equipo tiene
función Hz/Duty\,\%, verificar la frecuencia PWM (1\,kHz) sin osciloscopio.

\paragraph{Con simulación.} Construir cada caso completo (transistor + carga + diodo) en
NI Multisim/Proteus y correr un análisis transitorio: graficar $V_{CE}$, $I_C$, $I_B$ y
la corriente del diodo. La simulación permite además \textbf{observar de forma segura} lo
que ocurre \emph{sin} el diodo \emph{flyback} (pico de voltaje destructivo) — un
experimento que \textbf{nunca} debe hacerse con el circuito físico real. Es la forma
recomendada de estudiar la \S8.5 (nota de transitorio) cuando no se tiene acceso a un
osciloscopio de banco.

\paragraph{Requiere laboratorio presencial.} La forma de onda real del transitorio de
apagado (con y sin diodo) en la escala de nanosegundos a microsegundos, y el ruido de
conmutación real sobre el bus de 3.3\,V, solo pueden observarse con un osciloscopio
físico — la simulación es idealizada y no reproduce parásitos reales de cableado.

\begin{table}[H]\centering
\caption{Alternativas de trabajo en casa por caso de estudio.}
\begin{tabular}{lccc}
\toprule
Caso & Multímetro & Simulación & Requiere laboratorio \\
\midrule
1 (buzzer, señal) & Sí ($V_{CE(\text{sat})}$) & Sí & No \\
2 (motor, \emph{low-side}) & Sí & Sí & Recomendado (disipador, transitorio) \\
3 (motor, \emph{high-side}) & Sí & Sí & Recomendado (disipador, transitorio) \\
4 (ARINC 429 discreto) & --- & Sí (lógica) & No \\
Transitorio sin \emph{flyback} & No (peligroso) & \textbf{Sí, solo así} & --- \\
\bottomrule
\end{tabular}
\end{table}

% ============================================================================
\section{Herramientas de Verificación}
\subsection{Verificación con osciloscopio}
\begin{enumerate}
  \item \textbf{Señal PWM en la base}: verificar frecuencia (1\,kHz), amplitud (3.3\,V)
        y ciclo de trabajo.
  \item $V_{CE(\text{sat})}$ \textbf{durante saturación}: conectar CH1 al colector, CH2 al
        emisor. La diferencia debe ser $< V_{CE(\text{sat}),\max}$ del \emph{datasheet}.
  \item \textbf{Transitorio al apagado}: capturar el voltaje en el colector en modo
        \emph{single-shot} al pasar \emph{duty} de 100\,\% a 0\,\%. Con diodo
        \emph{flyback}: limitado a $\sim V_{CC}+0.3$\,V. Sin diodo: pico potencialmente
        destructivo ($>\SI{100}{\volt}$).
  \item \textbf{Corriente de base}: medir voltaje a ambos lados de $R_B$ y calcular
        $I_B=\Delta V/R_B$.
\end{enumerate}

% ============================================================================
\section{Registro de Datos y Análisis}

\subsection{Tabla de caracterización de los interruptores}
\begin{table}[H]\centering
\caption{Bitácora de caracterización de los interruptores electrónicos.}
\begin{tabular}{lcccc}
\toprule
Carga / Caso & $V_\text{GPIO}$ (V) & $I_B$ (mA) & $I_C$ (mA) & $V_{CE(\text{sat})}$ (V) \\
\midrule
Buzzer (Caso 1) & 3.3 & & & \\
Motor 12\,V (Caso 2) & 3.3 & & & \\
Motor 12\,V (Caso 3) & 3.3 & & & \\
Falla (sin \emph{flyback}) & 3.3 & & & Estado térmico: Crítico \\
\bottomrule
\end{tabular}
\end{table}

\subsection{Tabla de verificación de cálculos}
\begin{table}[H]\centering
\caption{Comparación entre valores calculados y medidos.}
\begin{tabular}{clccc}
\toprule
Caso & Parámetro & Calculado & Medido & Error (\%) \\
\midrule
1 & $I_B$ (mA) & 2.0 & & \\
1 & $V_{CE(\text{sat})}$ (V) & $\le 0.2$ & & \\
2 & $I_{B,Q2}$ (mA) & $\le 450$ & & \\
2 & $V_{CE(\text{sat}),Q2}$ (V) & $\le 1.4$ & & \\
2 & $P_{D,Q2}$ (W) & 2.1 & & \\
3 & $I_{B,Q2}$ (mA) & $\le 300$ & & \\
3 & $|V_{EC(\text{sat}),Q2}|$ (V) & $\le 1.4$ & & \\
\bottomrule
\end{tabular}
\end{table}

\subsection{Tabla de transitorios (osciloscopio)}
\begin{table}[H]\centering
\caption{Voltaje pico en el colector al apagar la carga.}
\begin{tabular}{ll}
\toprule
Condición & $V_\text{pico}$ (V) \\
\midrule
Con diodo \emph{flyback} & Esperado: $\sim 12.3$ \\
Sin diodo \emph{flyback} & Esperado: $> 50$ (solo en simulación, \S9) \\
\bottomrule
\end{tabular}
\end{table}

\subsection{Gráficas requeridas}
\begin{enumerate}
  \item $V_{CE(\text{sat})}$ vs $I_C$: para cada transistor, comparar punto de operación
        medido con la curva del \emph{datasheet}.
  \item Transitorio de apagado (osciloscopio o simulación): forma de onda en el colector
        con y sin diodo \emph{flyback}.
  \item PWM \emph{duty} vs potencia media: verificar relación lineal $P\propto D$ en el
        motor.
  \item Capturas de simulación Multisim: punto de operación en saturación para cada caso.
\end{enumerate}

% ============================================================================
\section{Análisis de Resultados}
\begin{enumerate}
  \item \textbf{Verificación de saturación}: comparar $V_{CE(\text{sat})}$ medido con el
        valor del \emph{datasheet}. Calcular la potencia disipada real
        $P=V_{CE}\times I_C$. Determinar si el transistor requiere disipador según
        $R_{\theta JA}$ y $T_{J,\max}$.
  \item \textbf{Análisis de transitorios}: si se usó osciloscopio o simulación, comparar
        la amplitud del pico con y sin diodo \emph{flyback}. Explicar el riesgo para el
        microcontrolador y el transistor.
  \item \textbf{Comparativa \emph{low-side} vs \emph{high-side}}: comparar complejidad,
        número de componentes, eficiencia y ventajas de seguridad de cada configuración.
  \item \textbf{Potencia en resistencias}: verificar que las resistencias de base y
        \emph{pull-up} no se sobrecalienten. Confirmar que la selección de potencia
        (1/4\,W vs 5\,W) fue correcta.
  \item \textbf{Integridad de control}: discutir la fiabilidad del control por bus
        ARINC 429. ¿Qué sucedería si el bit de la bomba cambiara por EMI sin validación
        de paridad?
  \item \textbf{Análisis de simulación vs experimento}: comparar los resultados de
        Multisim con las mediciones de laboratorio (o de casa, \S8). Identificar
        discrepancias y posibles causas (tolerancias, modelado).
\end{enumerate}

% ============================================================================
\section{Preguntas de Discusión y Refuerzo}
\begin{enumerate}
  \item \textbf{High-Side vs Low-Side}: ¿por qué en aeronáutica se prefiere que los
        interruptores de potencia corten la línea positiva (\emph{high-side}) en lugar
        de la línea de tierra? Considere la seguridad ante fallas de aislamiento a la
        estructura metálica del avión.
  \item \textbf{BJT vs MOSFET}: muchos sistemas modernos usan MOSFETs de potencia.
        Explique la ventaja en términos de control por voltaje (vs corriente),
        impedancia de entrada y $R_{DS(\text{on})}$ para eficiencia. ¿Por qué un MOSFET
        de canal N con $V_{GS(\text{th})}<\SI{2}{\volt}$ puede reemplazar al 2N2222A +
        TIP31C con un solo componente?
  \item \textbf{Factor de seguridad}: ¿por qué se calcula $I_B$ con ODF$>1$? Relacione
        con la degradación de $h_{FE}$ por temperatura y envejecimiento.
  \item \textbf{EMI conductiva}: el motor DC genera ruido eléctrico en el bus de 12\,V.
        ¿Cómo protegería la línea de 3.3\,V del MCU para evitar reinicios? Considere
        condensadores de desacoplo, ferrita, y separación física.
  \item Si la resistencia $R_\text{PU}$ del Caso 2 se calculó para \SI{22}{\ohm} a
        \SI{4.5}{\watt}, ¿qué ocurre si se usa una resistencia de 1/4\,W? Describa el
        modo de falla.
  \item ¿Cuál es el efecto de omitir el diodo \emph{flyback} en un solenoide de tren de
        aterrizaje que consume 2\,A a 28\,VDC? Estime el voltaje transitorio con la
        inductancia del solenoide ($L\approx\SI{50}{\milli\henry}$).
  \item En un sistema FBW con triple redundancia, ¿cómo diseñaría la etapa de potencia
        para un actuador de alerón? Considere: 3 canales independientes, votación,
        monitoreo cruzado y modo de reversión mecánica.
  \item Compare la eficiencia de conmutación de un BJT ($V_{CE(\text{sat})}=\SI{1.4}{\volt}$)
        con un MOSFET ($R_{DS(\text{on})}=\SI{0.05}{\ohm}$) para una carga de 1.5\,A.
        Calcule la potencia disipada en cada caso.
  \item ¿Qué norma (DO-160G, sección 22) especifica la susceptibilidad a transitorios
        inducidos por rayo en aviónica? ¿Cómo se relaciona esto con la protección
        \emph{flyback}?
\end{enumerate}

% ============================================================================
\section{Entregables (Formato AIAA)}
\begin{enumerate}
  \item \textbf{Código fuente}: \code{main.py} + \code{lib/actuator\_pwm.py}
        (MicroPython) o \code{p6.cpp} (C++).
  \item \textbf{Memoria de cálculo}: para los 3 casos — $I_B$, $R_B$, $P_D$, $T_J$,
        selección de diodo.
  \item \textbf{Capturas de simulación Multisim}: esquema y gráficas de análisis
        transitorio por caso.
  \item \textbf{Tabla de caracterización} (Registro \S10.1): $V_{CE(\text{sat})}$,
        $I_B$, $I_C$ medidos.
  \item \textbf{Verificación cálculos vs mediciones} (\S10.2): error porcentual.
  \item \textbf{Gráficas de transitorios}: con y sin diodo \emph{flyback} (osciloscopio
        o simulación).
  \item \textbf{Análisis}: mínimo 4 de los 6 puntos de la \S11.
  \item \textbf{Discusión}: responder $\ge 4$ de las 9 preguntas.
  \item \textbf{Conclusiones}: hallazgos, fuentes de error, mejoras propuestas.
  \item \textbf{Referencias}: $\ge 3$ fuentes (\emph{datasheets} de transistores y
        diodos, libro de electrónica, normativa DO-160G).
\end{enumerate}

% ============================================================================
\section{Referencias}
\begin{enumerate}
  \item Boylestad, R.L., Nashelsky, L. \emph{Electrónica: Teoría de Circuitos y
        Dispositivos Electrónicos}, 11a ed., Pearson, 2009.
  \item ON Semiconductor. \emph{2N2222A Datasheet}, Rev.\ 11, 2013.
  \item STMicroelectronics. \emph{TIP31C Datasheet}, Rev.\ 7.
  \item STMicroelectronics. \emph{TIP32C Datasheet}, Rev.\ 5.
  \item ON Semiconductor. \emph{1N5817--1N5819 Datasheet} (Schottky Barrier Rectifier).
  \item ARINC. \emph{Specification 429 Part 1}: Mark 33 Digital Information Transfer
        System (DITS), 2004.
  \item RTCA. \emph{DO-160G}: Environmental Conditions and Test Procedures for Airborne
        Equipment, Sec.\ 20 (RF Susceptibility), Sec.\ 22 (Lightning Induced Transient
        Susceptibility), 2010.
  \item RTCA. \emph{DO-178C}: Software Considerations in Airborne Systems and Equipment
        Certification, 2011.
  \item Moir, I., Seabridge, A. \emph{Aircraft Systems: Mechanical, Electrical, and
        Avionics Subsystems Integration}, 3rd ed., Wiley, 2008.
  \item MicroPython. \emph{machine.PWM Documentation}.
\end{enumerate}

% ============================================================================
\appendix
\section{Código Fuente del Repositorio}
El código completo está en \url{https://github.com/EU97/SISELA-Init}.
\begin{table}[H]\centering
\caption{Estructura de archivos de la Práctica 6.}
\begin{tabular}{ll}
\toprule
Plataforma & Archivo(s) / Descripción \\
\midrule
MicroPython/ESP32 & \code{P6/main.py} (4 modos), \code{P6/boot.py},
                    \code{P6/lib/actuator\_pwm.py} \\
MicroPython/RP2040 & \code{P6/main.py} (adaptación ADC 16-bit), \code{P6/boot.py},
                     \code{P6/PINES.md} \\
C++ / PlatformIO & \code{practices/p6.cpp}, \code{common/propulsion.h}
                   (\code{PropulsionSystem}, API 0--100\,\%) \\
\bottomrule
\end{tabular}
\end{table}

\subsection{Diferencias entre plataformas}
\begin{table}[H]\centering
\caption{Resumen de diferencias P6 entre plataformas.}
\begin{tabular}{lll}
\toprule
Aspecto & ESP32 & RP2040 \\
\midrule
Pin PWM & GPIO18 & GP18 (PWM1A) \\
Frecuencia PWM & 1\,kHz (LEDC) & 1\,kHz (\emph{slice} PWM) \\
Jitter PWM & $\sim\SI{10}{\nano\second}$ & $\sim\SI{1}{\nano\second}$ \\
Pin ADC & GPIO34 (solo entrada) & GP26 (ADC0) \\
Resolución ADC & 12 bits (0--4095) & 16 bits (0--65535) \\
Lectura ADC & \code{adc.read()} & \code{adc.read\_u16()} \\
Config.\ ADC & \code{atten(11dB)} & Ninguna \\
Alimentación 5\,V & Fuente ext.\ siempre & VSYS para cargas $<\SI{500}{\milli\ampere}$ \\
C++ PWM & \code{analogWrite} & \code{analogWrite} (core Arduino) \\
\bottomrule
\end{tabular}
\end{table}

\section{Contexto Aeronáutico Complementario}
\subsection{Sistemas de potencia en aeronaves}
En aeronaves modernas, los buses de potencia eléctrica operan típicamente a 28\,VDC
(aviación general y aviones militares) o 115\,VAC / 270\,VDC (transporte comercial con
sistemas más eléctricos). Los interruptores de estado sólido (SSPCs — \emph{Solid State
Power Controllers}) han reemplazado progresivamente a los \emph{circuit breakers}
mecánicos, ofreciendo protección programable, diagnóstico remoto y menor peso. Los SSPCs
utilizan MOSFETs de potencia en configuración \emph{high-side} con circuitos de
protección integrados: sobrecorriente (fusible electrónico programable), arco eléctrico
(\emph{arc fault}, requisito FAR 25.1353), monitoreo de corriente para diagnóstico, y
comunicación por bus ARINC 429, ARINC 629 o CAN aeroespacial.

\subsection{Conmutación high-side en aviónica}
La preferencia por conmutación \emph{high-side} en aeronaves se debe a la práctica
estándar de conexión a tierra (\emph{grounding}). En la mayoría de las aeronaves, la
estructura metálica del fuselaje sirve como retorno de corriente (negativo del bus DC).
Si un cable de carga toca la estructura: con interruptor \emph{high-side}, la carga está
desconectada del positivo $\Rightarrow$ no hay cortocircuito; con interruptor
\emph{low-side}, la carga está permanentemente conectada al positivo $\Rightarrow$
cortocircuito a tierra $\Rightarrow$ falla potencialmente catastrófica.

\end{document}
